\scanpage{88}
Par définition $=$ de $\Omega^+$ etc, il est clair que les images sont les $\Omega^+$ etc (mais a-t-on $\Omega^+=\Omega$ etc ?) \margin{?} L'image de $\unsure{E|\mathcal E}$\inserted{$=(\gamma\unsure{*|\times}\gamma)^{(*{\unsure{G|\mathfrak S|\mathcal S}}_3/\pi_{01})}$} resp. $E_0$ est évidemment $|E|$ (resp. $|E|\unsure{.|{}_0}$) \ldots

On a un hom d'inclusion $\gamma_0\mathbin{-}\gamma$, définissant
\diagram{\begin{tikzcd}[ampersand replacement=\&,row sep=tiny]
(\gamma_0\times\gamma_0)^{*(G_3/\pi_{01})}\arrow[r] \&(\gamma*\gamma)^{(*G_3/\pi_{01})}\\
\rotatebox{90}{$=$}\&\rotatebox{90}{$=$}\\
E_0\&E
\end{tikzcd}}
et donnant lieu à diagramme comm.
\diagram{\begin{tikzcd}[ampersand replacement=\&,column sep=large]
G_3\underset{1/2}{\cdot} E_0\arrow[r,dash]\arrow[d] \& G_3\cdot E\arrow[d]\\
\Omega_0\arrow[r] \& \Omega.
\end{tikzcd}}

On a aussi, notant que
\diagram{\begin{tikzcd}[ampersand replacement=\&,column sep=small]
\gamma \&[-1em] = \&[-1em] \overset{\substack{\deleted{\unsure{\gamma}}\\\rotatebox{90}{$=$}}}{\widehat{\mathbb Z}^{*}}\arrow[r,dash,"{\deleted{\unsure{\sim}}}"] \&(\mathbb Z/6\mathbb Z)^*\arrow[r,"\sim"]\&(\mathbb Z/3\mathbb Z)^*\arrow[r,dash,"\sim"]\&\pm1\\
\gamma_0\arrow[u,hook] \&[-1em] = \&[-1em] \mathbb Z^*\arrow[u,hook]\arrow[urrr,bend right=16,"\sim"']\&\deleted{\unsure{\gamma_2|\chi_2}}\ \deleted{\unsure{\gamma_0=\mathbb Z^*|\chi_0=\mathbb Z^*}}\arrow[u]\&\&
\end{tikzcd}}
un hom $\gamma\longrightarrow\gamma_0$ inv. à g. de $\gamma_0\hookrightarrow\gamma$, d'où $E\longrightarrow E_0$ inv à g. de $E_0\mathbin{-}E$, d'où
\[
G_3\cdot E\mathbin{\text{---}}G_3\cdot E_0
\]
inv. à g. de \unsure{tt} à l'heure. D'autre part considérons
\[
\deleted{\unsure{\Omega}}\Pi(6)=\text{quotient de }\Pi\text{ par \deleted{la} \inserted{\unsure{(l'introduct de)|(l'adhérence de)}} relation }
\underset{\deleted{\unsure{h^6}}\ g^6=1\quad(g\in\Pi_0)}{\vphantom{\Pi}\strut}
\]
% archive sheet number: 87
